\documentclass[12pt]{article}
\usepackage[margin=1in]{geometry}
\usepackage{titlesec}
\usepackage{hyperref}
\usepackage{booktabs}
\usepackage{longtable}
\usepackage{array}
\usepackage{parskip}
\usepackage{enumitem}
\usepackage{fancyhdr}

\hypersetup{colorlinks=true, linkcolor=blue, urlcolor=blue}

\pagestyle{fancy}
\fancyhf{}
\rhead{Campus Resource Booking \& Check-In System}
\lhead{SRS}
\rfoot{Page \thepage}

\title{
    \vspace{2cm}
    {\Large Software Requirements Specification (SRS)} \\[1em]
    {\LARGE \textbf{Campus Resource Booking \& Check-In System}} \\[2em]
    {\large CS3338 --- Group 1} \\[0.5em]
    {\normalsize Version 4.0}
    \vspace{2cm}
}
\author{Herlyn Dsouza}
\date{May 15, 2025}

\begin{document}

\maketitle
\thispagestyle{empty}
\newpage

\tableofcontents
\newpage

% -------------------------------------------------------
\section*{Versions Table}
\addcontentsline{toc}{section}{Versions Table}

\begin{longtable}{|p{1.4cm}|p{2.4cm}|p{3.5cm}|p{6.5cm}|}
\hline
\textbf{Ver.} & \textbf{Date} & \textbf{Author(s)} & \textbf{Changes} \\
\hline
1.0 & Feb 2, 2025  & Herlyn Dsouza & Initial draft. Introduction, user roles, core interface requirements, legal and ethical section. \\
\hline
2.0 & Mar 2, 2025  & Herlyn Dsouza & Added full REST API endpoint table, auth requirements, booking creation and cancellation requirements, input validation rules. \\
\hline
3.0 & Apr 6, 2025  & Herlyn Dsouza & Added check-in requirements (time window), admin dashboard requirements, analytics endpoint spec, no-show handling requirement. \\
\hline
4.0 & May 15, 2025 & Herlyn Dsouza & Final revision. Added rate limiting requirement, session/token expiry requirements, performance requirements, finalized all sections. \\
\hline
\end{longtable}
\newpage

% -------------------------------------------------------
\section{Introduction}

\subsection{Purpose}
This Software Requirements Specification (SRS) defines the complete functional and non-functional requirements for the Campus Resource Booking \& Check-In System. It provides a shared understanding between the development team, QA, and the course instructor regarding what the system must do and the constraints under which it must operate.

\subsection{Intended Audience}
\begin{itemize}
    \item \textbf{Development team} --- Britany Rodriguez, Herlyn Dsouza, Israel Melchor (to guide implementation decisions at each snapshot)
    \item \textbf{DevOps / QA} --- Nautica Uribe (to design TestRail test cases that verify each requirement)
    \item \textbf{Course Instructor} --- for evaluation and project approval
\end{itemize}

\subsection{Project Overview}
The Campus Resource Booking \& Check-In System is a web-based platform enabling students to reserve campus resources (study rooms, labs, tutoring appointments) and physically check in upon arrival. Staff manage their resources; admins oversee the entire platform and view usage analytics.

\subsection{Definitions}
\begin{itemize}
    \item \textbf{Resource} --- A bookable campus asset: study room, computer lab seat, or tutoring appointment slot
    \item \textbf{Booking} --- A reservation made by a student for a specific resource and time slot
    \item \textbf{Check-In} --- The action a student performs on arrival to confirm attendance at a booked resource
    \item \textbf{Conflict} --- A state where two active bookings overlap in time for the same resource
    \item \textbf{No-Show} --- A booking whose check-in window expired without the student checking in
    \item \textbf{Check-In Window} --- The period during which a student may check in: 15 min before to 10 min after booking start
\end{itemize}

% -------------------------------------------------------
\section{External Interface Requirements}

\subsection{User Interface Requirements}

\subsubsection{General UI --- All Snapshots}
\begin{itemize}
    \item The system shall provide a responsive web interface accessible from modern desktop and mobile browsers (Chrome, Firefox, Safari, Edge).
    \item All primary pages shall render and be interactive within 3 seconds on a standard broadband connection.
    \item The UI shall use a consistent navigation bar so users can reach any core feature within 2 clicks.
    \item All forms shall display inline validation errors before submission.
    \item The application shall display loading indicators during all API requests.
\end{itemize}

\subsubsection{Authentication Interface --- Snapshot 2}
\begin{itemize}
    \item The system shall provide a \textbf{Login} page with email and password fields plus a submit button.
    \item The system shall provide a \textbf{Register} page collecting full name, email, password, and role (Student/Staff).
    \item Passwords shall always be masked in input fields.
    \item Invalid login attempts shall display a generic ``Invalid credentials'' message without revealing whether the email exists.
    \item Successful login shall redirect the user to the Resource Browse page.
    \item Successful registration shall redirect the user to the Login page with a confirmation message.
\end{itemize}

\subsubsection{Student Interface --- Snapshot 2}
\begin{itemize}
    \item Students shall be able to browse all active resources with filtering by type and location.
    \item Students shall be able to select a date and view available time slots for any resource.
    \item Students shall be able to create a booking by selecting a resource and an available time slot.
    \item Students shall be able to view all their upcoming and past bookings.
    \item Students shall be able to cancel any upcoming booking that is not within the check-in window.
\end{itemize}

\subsubsection{Check-In Interface --- Snapshot 3}
\begin{itemize}
    \item The Check-In button shall appear on the My Bookings page for any upcoming booking.
    \item The Check-In button shall be \textbf{active} (clickable) only within the check-in window (15 min before to 10 min after start time).
    \item Outside the check-in window, the button shall be disabled with a tooltip explaining why.
    \item Successful check-in shall update the booking status badge to ``Checked In'' without requiring a page reload.
\end{itemize}

\subsubsection{Staff Interface --- Snapshots 2 \& 3}
\begin{itemize}
    \item Staff shall be able to create new resources with name, type, location, capacity, and available hours.
    \item Staff shall be able to edit and deactivate resources they manage.
    \item Staff shall be able to view a list of all bookings for their resources.
\end{itemize}

\subsubsection{Admin Interface --- Snapshot 3}
\begin{itemize}
    \item Admins shall have a dedicated dashboard accessible only to the \texttt{admin} role.
    \item The dashboard shall display: total bookings, active resources, check-in rate, and no-show rate as summary cards.
    \item The dashboard shall display a bar chart of bookings per day for the last 30 days.
    \item Admins shall be able to manage (view, edit role, deactivate) all users.
    \item Admins shall be able to view and cancel any booking on the platform.
\end{itemize}

\subsection{Software Interfaces}

\subsubsection{Frontend--Backend Interface}
\begin{itemize}
    \item All frontend-backend communication shall use REST over HTTP/HTTPS.
    \item All authenticated requests shall include a valid JWT in the \texttt{Authorization: Bearer} header.
    \item The API shall return JSON responses with correct HTTP status codes (200, 201, 400, 401, 403, 404, 409, 429, 500).
    \item Validation errors shall return HTTP 400 with a JSON array of field-level error messages.
\end{itemize}

\subsubsection{Backend--Database Interface}
\begin{itemize}
    \item The backend shall use Mongoose to communicate with MongoDB.
    \item All database operations shall be wrapped in \texttt{try/catch} blocks.
    \item Booking conflict checks shall be performed as a server-side MongoDB query before any booking is confirmed.
\end{itemize}

\subsubsection{REST API Endpoint Specification --- Snapshot 2}

\begin{longtable}{|p{1.4cm}|p{4.5cm}|p{2cm}|p{5cm}|}
\hline
\textbf{Method} & \textbf{Endpoint} & \textbf{Auth} & \textbf{Description} \\
\hline
POST & /api/auth/register & Public & Register a new user \\
\hline
POST & /api/auth/login    & Public & Login, receive JWT \\
\hline
GET  & /api/resources     & Any    & List all active resources \\
\hline
POST & /api/resources     & Staff/Admin & Create a resource \\
\hline
GET  & /api/resources/:id & Any    & Get resource details \\
\hline
PUT  & /api/resources/:id & Staff/Admin & Update a resource \\
\hline
DELETE & /api/resources/:id & Admin  & Deactivate a resource \\
\hline
GET  & /api/bookings      & Student+ & Get user's bookings \\
\hline
POST & /api/bookings      & Student+ & Create a booking \\
\hline
PUT  & /api/bookings/:id/cancel  & Student+ & Cancel a booking \\
\hline
\end{longtable}

\subsubsection{Additional Endpoints --- Snapshot 3}

\begin{longtable}{|p{1.4cm}|p{4.5cm}|p{2cm}|p{5cm}|}
\hline
\textbf{Method} & \textbf{Endpoint} & \textbf{Auth} & \textbf{Description} \\
\hline
PUT & /api/bookings/:id/checkin & Student & Check in to booking (window-validated) \\
\hline
GET & /api/admin/users & Admin & List all users \\
\hline
PUT & /api/admin/users/:id & Admin & Update user role or status \\
\hline
GET & /api/admin/bookings & Admin & List all bookings \\
\hline
GET & /api/admin/analytics & Admin & Usage analytics data \\
\hline
\end{longtable}

\subsection{Non-Functional Requirements --- Snapshot 4}
\begin{itemize}
    \item \textbf{Rate Limiting:} Authentication endpoints (\texttt{/api/auth/*}) shall enforce a maximum of 10 requests per IP address per 15-minute window. Exceeding this limit shall return HTTP 429.
    \item \textbf{Token Expiry:} All JWTs shall expire after 24 hours. Expired tokens shall be rejected with HTTP 401.
    \item \textbf{Password Security:} All passwords shall be hashed with bcrypt at a minimum cost factor of 10.
    \item \textbf{HTTPS:} In production deployment, all traffic shall be served over HTTPS only.
    \item \textbf{Performance:} API endpoints shall respond within 500ms under normal load conditions.
\end{itemize}

% -------------------------------------------------------
\section{Legal and Ethical Considerations}

\subsection{User Data and Privacy}
\begin{itemize}
    \item The system collects personally identifiable information (PII): names, email addresses, and booking history.
    \item Passwords shall never be stored in plaintext or logged at any level.
    \item All user data shall remain within the system and shall not be shared with any third party.
    \item The system adheres to the principle of \textbf{data minimization}: only the data necessary to operate the core features is collected.
    \item Users shall be informed of data collection upon registration.
\end{itemize}

\subsection{Access Control}
\begin{itemize}
    \item Every API endpoint shall enforce role-based access control (RBAC).
    \item A student shall not be able to call staff or admin routes, even with a valid JWT.
    \item A staff member shall only be able to view/edit resources they manage.
    \item JWT tokens expire after 24 hours, requiring re-authentication.
\end{itemize}

\subsection{Data Retention}
\begin{itemize}
    \item Booking records shall be retained for a minimum of one academic semester to support analytics and dispute resolution.
    \item Users shall be able to request account deletion, which shall remove or anonymize their PII from the database.
\end{itemize}

\subsection{Ethical Considerations}
\begin{itemize}
    \item \textbf{Equitable access:} The booking system is first-come, first-served. No student group receives preferential treatment.
    \item \textbf{No surveillance:} Check-in data is used solely to confirm attendance for the booked resource. It shall not be used to build behavioral profiles or track movements.
    \item \textbf{Transparency:} Users are clearly informed what data is collected and why.
    \item \textbf{Minimal privilege:} Staff accounts can only manage their assigned resources and cannot view other users' personal profiles.
    \item \textbf{Disability access:} The UI shall follow WCAG 2.1 AA accessibility guidelines where feasible, including keyboard navigation and sufficient color contrast.
\end{itemize}

% -------------------------------------------------------
\section*{Glossary}
\addcontentsline{toc}{section}{Glossary}

\begin{longtable}{|p{3cm}|p{10cm}|}
\hline
\textbf{Term / Acronym} & \textbf{Definition} \\
\hline
SRS   & Software Requirements Specification \\
\hline
API   & Application Programming Interface \\
\hline
REST  & Representational State Transfer \\
\hline
JWT   & JSON Web Token \\
\hline
RBAC  & Role-Based Access Control \\
\hline
PII   & Personally Identifiable Information \\
\hline
CRUD  & Create, Read, Update, Delete \\
\hline
MVP   & Minimum Viable Product \\
\hline
ODM   & Object Document Mapper \\
\hline
HTTPS & Hypertext Transfer Protocol Secure \\
\hline
WCAG  & Web Content Accessibility Guidelines \\
\hline
\end{longtable}

\end{document}
